% DRAFT of 17 November 2026, author-side session. Fragment in the notation of Appendix G of
% "From Correlations to Geometry" (Lemma lem:consistency, Theorem thm:kernel).
% STATUS: the proposition is a deterministic statement on the event of Lemma lem:consistency and is
% proven here; it has NOT yet been read by the second session. It becomes useful only together with
% a bound for the statistic Pi_k (being written by the second session). Not read by a person.

\begin{proposition}[Reduction of the trace to a polynomial statistic]\label{prop:reduction}
In the situation of Theorem~\ref{thm:kernel} let
\[
\theta^1_{bi}=-\epsilon_i-\epsilon_b+(\tilde P\epsilon)_b,\qquad
\hat h^{(i)}_m(b)=\sum_j\tilde P_{bj}\,(1-\epsilon_j)\,\big(p_m(x_j-x_i)-p_{m+1}(x_b-x_i)\big),
\]
and, for $k\ge1$,
\[
\Pi_k=\frac1{N^2}\sum_{i,b}p_k(x_i-x_b)\,p(x_b-x_i)\,\theta^1_{bi}
+\sum_{m=1}^{k-1}\frac1{N^2}\sum_{i,b}p_{k-m}(x_i-x_b)\,\hat h^{(i)}_m(b)
+\frac1N\sum_i\hat h^{(i)}_k(i).
\]
Let $\varphi_d(k)=1$ for $d\le2$, $\varphi_3(k)=1+\log k$ and $\varphi_d(k)=k^{(d-3)/2}$ for $d\ge4$. There is a constant $C_2$ depending only on $d$ such that on the event of Lemma~\ref{lem:consistency}, for all integers $1\le k<(C_1\delta')^{-2}$,
\[
\big|\Tr P^{k+1}-m_{k+1}-\Pi_k\big|\le C_2\,\delta'^2\,k\,\varphi_d(k)\,m_{k+1}.
\]
$\Pi_k$ is a polynomial of degree at most four in the values of the kernels $p_s(x_a-x_b)$.
\end{proposition}

\begin{proof}
$C$ denotes constants that depend only on $d$. We use the notation of the proofs of Lemma~\ref{lem:consistency} and Theorem~\ref{thm:kernel}; in particular $\delta\le6\delta'\le\frac18$, $u_j=(1+\epsilon_j)^{-1}$, $r_b=\sum_j\tilde P_{bj}u_j$, and the proof of the lemma shows that on its event, besides (ii), also $A\le C\,m^{-1/2}p_{6m}(x_b-x_i)$ for the quantity $A=\frac1N\sum_jp(x_b-x_j)|p_m(x_j-x_i)-p_{m+1}(x_b-x_i)|$.

\emph{Step 1: an identity.} By Step~2 of the proof of Theorem~\ref{thm:kernel}, $N(P^{k+1})_{ii}=(P^kRg_i)(i)+\sum_b(P^k)_{ib}\,p(x_b-x_i)\,\theta_{bi}$, and by \eqref{eq:kernel-telescope} with $a=i$,
\[
\Tr P^{k+1}-m_{k+1}=\frac1N\sum_i\sum_b(P^k)_{ib}\,p(x_b-x_i)\,\theta_{bi}+\frac1N\sum_i\sum_{m=1}^k\big(P^{k-m}h^{(i)}_m\big)(i).
\]
\emph{Step 2: linearisation in $\epsilon$.} Since $\sum_j\tilde P_{bj}=1+\epsilon_b$ and $u_j\in[\frac1{1+\delta},\frac1{1-\delta}]$, $|r_b-1|\le\frac{2\delta}{1-\delta}\le\frac{16}7\delta$. Write $u_j=1-\epsilon_j+\vartheta_j$ with $0\le\vartheta_j=\epsilon_j^2/(1+\epsilon_j)\le\frac87\delta^2$; then $r_b-1=\epsilon_b-(\tilde P\epsilon)_b+(\tilde P\vartheta)_b$ and $|(\tilde P\vartheta)_b|\le\frac97\delta^2$. With $\varrho=r_b-1$ and $(1+\varrho)^{-1}=1-\varrho+\varrho^2/(1+\varrho)$,
\[
\theta_{bi}=\frac{u_i}{r_b}-1=\theta^1_{bi}-(\tilde P\vartheta)_b+\epsilon_i\varrho+\vartheta_i(1-\varrho)+u_i\frac{\varrho^2}{1+\varrho},\qquad\text{hence}\qquad|\theta_{bi}-\theta^1_{bi}|\le14\,\delta^2 .
\]
For the defect, \eqref{eq:cons-split} gives $r_b\,h^{(i)}_m(b)=S_1+S_2$ and $S_2=-\sum_j\tilde P_{bj}\epsilon_j(g(x_j)-c_b)+\sum_j\tilde P_{bj}\vartheta_j(g(x_j)-c_b)$, so that $S_1+S_2=\hat h^{(i)}_m(b)+E'$ with $|E'|\le\frac87\delta^2A$, and
\[
h^{(i)}_m(b)-\hat h^{(i)}_m(b)=(1-r_b)\,h^{(i)}_m(b)+E',\qquad\big|h^{(i)}_m(b)-\hat h^{(i)}_m(b)\big|\le C\,\delta'^2\,m^{-1/2}\,p_{6m}(x_b-x_i).
\]
With (ii) and (iii) of the lemma: $\max_b|h-\hat h|\le C\delta'^2m^{-1/2}m_{6m}$, $\frac1N\sum_b|h-\hat h|\le C\delta'^2m^{-1/2}$ and $\frac1N\sum_b|\hat h^{(i)}_m(b)|\le C\delta'm^{-1/2}$.

\emph{Step 3: the three kinds of terms.} (a) The last step of the walk. Since $|(P^k)_{ib}-p_k(x_i-x_b)/N|\le D_k/N$, $|\theta_{bi}|\le4\delta$, $p_k\le m_k$ and $\frac1N\sum_bp(x_b-x_i)\le1+\delta$,
\[
\Big|\frac1N\sum_{i,b}(P^k)_{ib}\,p(x_b-x_i)\,\theta_{bi}-\frac1{N^2}\sum_{i,b}p_k(x_i-x_b)\,p(x_b-x_i)\,\theta^1_{bi}\Big|\le(1+\delta)\big(4\delta\,D_k+14\,\delta^2m_k\big)\le C\,\delta'^2\sqrt k\;m_k .
\]
(b) For $1\le m\le k-1$ and $l=k-m\ge1$,
\[
\frac1N\sum_i\big(P^lh^{(i)}_m\big)(i)-\frac1{N^2}\sum_{i,b}p_l(x_i-x_b)\,\hat h^{(i)}_m(b)
=\frac1N\sum_{i,b}(P^l)_{ib}\,\big(h-\hat h\big)^{(i)}_m(b)+\frac1N\sum_{i,b}\Big[(P^l)_{ib}-\frac{p_l(x_i-x_b)}N\Big]\hat h^{(i)}_m(b).
\]
The first sum is at most $(m_l+D_l)\max_i\frac1N\sum_b|h-\hat h|\le C\delta'^2m^{-1/2}m_l$ in absolute value, the second at most $D_l\max_i\frac1N\sum_b|\hat h|\le C\delta'^2\sqrt l\,m^{-1/2}m_l$, by Theorem~\ref{thm:kernel}. (c) For $m=k$: $|(h-\hat h)^{(i)}_k(i)|\le C\delta'^2k^{-1/2}m_{6k}\le C\delta'^2m_k$.

\emph{Step 4: summation.} For $l\le k$, $m_l\le(k/l)^{d/2}m_k$, because $\gamma_a\le(t/a)^{d/2}\gamma_t$ for $a\le t$. Hence
\[
\sum_{l=1}^{k-1}\frac{\sqrt l\;m_l}{\sqrt{k-l}}\le m_k\,k^{d/2}\Big[\Big(\frac2k\Big)^{1/2}\sum_{l\le k/2}l^{(1-d)/2}+\Big(\frac2k\Big)^{(d-1)/2}\sum_{k/2<l<k}(k-l)^{-1/2}\Big]\le C\,k\,\varphi_d(k)\,m_k,
\]
since $\sum_{l\le k/2}l^{(1-d)/2}$ is at most $k/2$, $\sqrt{2k}$, $1+\log k$ and a constant for $d=1$, $d=2$, $d=3$ and $d\ge4$, and $\sum_{k/2<l<k}(k-l)^{-1/2}\le2\sqrt k$. Adding (a), (b) and (c) and using $m_k\le2^{d/2}m_{k+1}$ gives the claim. The degree of $\Pi_k$: $\hat h$ and $\theta^1$ are linear in $\epsilon$, and $\epsilon_j=\frac1N\sum_v(p(x_j-x_v)-1)$.
\end{proof}

% What this is for. If E|Pi_k| <= C (M/N) k m_k uniformly in k (lemma being written by the second session),
% then Tr P^{k+1}/m_{k+1} - 1 = O_P( (M log N / N) k phi_d(k) ), and the sharp rate xi^2/t holds in
% probability for N xi^(d+4)/log N -> infinity in dimensions d <= 2 (with a logarithm more for d = 3).
% Numerical check (one sample, d = 2, xi = 0.05, N = 1500 and 3000; script not in the repository):
% the remainder divided by delta'^2 k m_{k+1} is 0.12, 0.04, 0.02, 0.013, 0.007 for k = 1, 2, 4, 8, 16.
% The remainder itself is of the order (M/N) m_{k+1} and larger than Pi_k: Pi_k is NOT the leading part
% of the error, the second order in epsilon contributes at the same order in the mean.
